% GENERATED FILE. Edit canonical lessons, then run npm run generate:books.
% canonical-source-hash: fnv1a64:1998ebff2ae054a4
% canonical-lessons: TE-C147-tuducu, TE-C147-kattu, TE-C147-kuttu, TE-C147-gelucu, TE-C147-dacu

\chapter{To wipe, To tie, To sew, To win, To hide}
\label{ch:te-to-wipe-to-tie-to-sew-to-win-to-hide}
\hlchaptermodality{147}

\begin{chapteropening}
\textbf{By the end of this chapter:} \emph{I can say to wipe, to tie, to sew, to win and to hide.}

Complete the last lesson of chapter 147: I can say to wipe, to tie, to sew, to win and to hide.
\end{chapteropening}

\section[\texorpdfstring{\te{తుడుచు} (tuḍucu)}{తుడుచు (tuḍucu)}]{\te{తుడుచు} (tuḍucu) — to wipe}

\label{lesson:TE-C147-tuducu}



\begin{quote}
Before the new one: say the Telugu for to count, then the Telugu for to wash.
\end{quote}

\subsection*{You'll want to know: \te{తుడుచు}}
\textbf{\te{తుడుచు}} — \emph{tuḍucu} — \textquotedblleft{}to wipe\textquotedblright{}.

\begin{grammarlens}[title={What you've built}]
One more word for this chapter.
\end{grammarlens}

\subsection*{Your turn}
\begin{itemize}
\raggedright
  \item \emph{Say it:} \emph{tuḍucu}
  \item \emph{Say it:} \emph{tuḍucu}, once more
  \item \emph{Say it:} say \emph{tuḍucu}
  \item \emph{Recall:} say the Telugu for to count, then the Telugu for to wash, then say \emph{tuḍucu} again
\end{itemize}

\begin{tcolorbox}[breakable,colback=teal!4,colframe=teal!35!black,title={Before you move on}]
What does \te{తుడుచు} mean? (\textquotedblleft{}To wipe\textquotedblright{}.) Say it once more.
\end{tcolorbox}
\section[\texorpdfstring{\te{కట్టు} (kaṭṭu)}{కట్టు (kaṭṭu)}]{\te{కట్టు} (kaṭṭu) — to tie, to build}

\label{lesson:TE-C147-kattu}



\begin{quote}
Before the new one: say the Telugu for to wash, then the Telugu for to wipe.
\end{quote}

\subsection*{You'll want to know: \te{కట్టు}}
\textbf{\te{కట్టు}} — \emph{kaṭṭu} — \textquotedblleft{}to tie, to build\textquotedblright{}.

\begin{grammarlens}[title={What you've built}]
One more word for this chapter.
\end{grammarlens}

\subsection*{Your turn}
\begin{itemize}
\raggedright
  \item \emph{Say it:} \emph{kaṭṭu}
  \item \emph{Say it:} \emph{kaṭṭu}, once more
  \item \emph{Say it:} say \emph{kaṭṭu}
  \item \emph{Recall:} say the Telugu for to wash, then the Telugu for to wipe, then say \emph{kaṭṭu} again
\end{itemize}

\begin{tcolorbox}[breakable,colback=teal!4,colframe=teal!35!black,title={Before you move on}]
What does \te{కట్టు} mean? (\textquotedblleft{}To tie, to build\textquotedblright{}.) Say it once more.
\end{tcolorbox}
\section[\texorpdfstring{\te{కుట్టు} (kuṭṭu)}{కుట్టు (kuṭṭu)}]{\te{కుట్టు} (kuṭṭu) — to sew}

\label{lesson:TE-C147-kuttu}



\begin{quote}
Before the new one: say the Telugu for to wipe, then the Telugu for to tie.
\end{quote}

\subsection*{You'll want to know: \te{కుట్టు}}
\textbf{\te{కుట్టు}} — \emph{kuṭṭu} — \textquotedblleft{}to sew\textquotedblright{}.

\begin{grammarlens}[title={What you've built}]
One more word for this chapter.
\end{grammarlens}

\subsection*{Your turn}
\begin{itemize}
\raggedright
  \item \emph{Say it:} \emph{kuṭṭu}
  \item \emph{Say it:} \emph{kuṭṭu}, once more
  \item \emph{Say it:} say \emph{kuṭṭu}
  \item \emph{Recall:} say the Telugu for to wipe, then the Telugu for to tie, then say \emph{kuṭṭu} again
\end{itemize}

\begin{tcolorbox}[breakable,colback=teal!4,colframe=teal!35!black,title={Before you move on}]
What does \te{కుట్టు} mean? (\textquotedblleft{}To sew\textquotedblright{}.) Say it once more.
\end{tcolorbox}
\section[\texorpdfstring{\te{గెలుచు} (gelucu)}{గెలుచు (gelucu)}]{\te{గెలుచు} (gelucu) — to win}

\label{lesson:TE-C147-gelucu}



\begin{quote}
Before the new one: say the Telugu for to tie, then the Telugu for to sew.
\end{quote}

\subsection*{You'll want to know: \te{గెలుచు}}
\textbf{\te{గెలుచు}} — \emph{gelucu} — \textquotedblleft{}to win\textquotedblright{}.

\begin{grammarlens}[title={What you've built}]
One more word for this chapter.
\end{grammarlens}

\subsection*{Your turn}
\begin{itemize}
\raggedright
  \item \emph{Say it:} \emph{gelucu}
  \item \emph{Say it:} \emph{gelucu}, once more
  \item \emph{Say it:} say \emph{gelucu}
  \item \emph{Recall:} say the Telugu for to tie, then the Telugu for to sew, then say \emph{gelucu} again
\end{itemize}

\begin{tcolorbox}[breakable,colback=teal!4,colframe=teal!35!black,title={Before you move on}]
What does \te{గెలుచు} mean? (\textquotedblleft{}To win\textquotedblright{}.) Say it once more.
\end{tcolorbox}
\section[\texorpdfstring{\te{దాచు} (dācu)}{దాచు (dācu)}]{\te{దాచు} (dācu) — to hide}

\label{lesson:TE-C147-dacu}



\begin{quote}
Before the new one: say the Telugu for to sew, then the Telugu for to win.
\end{quote}

\subsection*{You'll want to know: \te{దాచు}}
\textbf{\te{దాచు}} — \emph{dācu} — \textquotedblleft{}to hide\textquotedblright{}.

\begin{grammarlens}[title={What you've built}]
One more word for this chapter.
\end{grammarlens}

\subsection*{Your turn}
\begin{itemize}
\raggedright
  \item \emph{Say it:} \emph{dācu}
  \item \emph{Say it:} \emph{dācu}, once more
  \item \emph{Say it:} say \emph{dācu}
  \item \emph{Recall:} say the Telugu for to sew, then the Telugu for to win, then say \emph{dācu} again
\end{itemize}

\begin{tcolorbox}[breakable,colback=teal!4,colframe=teal!35!black,title={Before you move on}]
What does \te{దాచు} mean? (\textquotedblleft{}To hide\textquotedblright{}.) Say it once more.
\end{tcolorbox}
